\begin{frame}{모델이 ``공짜''인 경우: 체스와 바둑}
  체스·바둑 같은 게임에는 특별한 사정이 있다 ---
  \begin{itemize}
    \item 규칙이 \textbf{완벽하게 알려져} 있다
    \item ``이 수를 두면 판이 정확히 이렇게 바뀐다''는 전이확률
      $P(s'|s,a)$를 Dyna-Q처럼 데이터로부터 배울 필요가
      \textbf{전혀 없음} (Week 4 모델 기반 방법의 이상적 조건)
    \item 문제는 \textbf{상태 수가 너무 많다}: 바둑판의 가능한
      국면 수는 약 $2 \times 10^{170}$으로 추정
  \end{itemize}
  \vskip0.6em
  \begin{columns}[T]
    \begin{column}{0.5\textwidth}
      \kb{Week 4.3 가치반복 같은 ``전체 계산''}
      \\[0.4em]
      {\footnotesize 모든 상태를 한 번씩 순회하며 계산 ---
      $10^{170}$ 개 상태에서는 \textbf{근본적으로 불가능}.}
    \end{column}
    \begin{column}{0.5\textwidth}
      \kb{MCTS(15.2)의 ``지금 근처만 깊이''}
      \\[0.4em]
      {\footnotesize ``지금 당장 놓인 상황 근처만'' 깊이 파고드는
      탐색이 필요해진다 --- 15.2절 MCTS가 풀려는 문제.}
    \end{column}
  \end{columns}
\end{frame}
